% GENERATED FILE. Edit canonical lessons, then run npm run generate:books.
% canonical-source-hash: fnv1a64:9ff321f812992f95
% canonical-lessons: KA-C381-pariganisu, KA-C381-bhavisu, KA-C381-viveka, KA-C381-manana-madu, KA-C381-buddhi

\chapter{Thinking It Through}
\label{ch:ka-thinking-it-through}
\hlchaptermodality{381}

\begin{chapteropening}
\textbf{By the end of this chapter:} \emph{I can say to consider, to suppose, good sense, to reflect on and intellect.}

Complete the last lesson of chapter 381: I can say to consider, to suppose, good sense, to reflect on and intellect.
\end{chapteropening}

\section[\texorpdfstring{\kn{ಪರಿಗಣಿಸು} (parigaṇisu)}{ಪರಿಗಣಿಸು (parigaṇisu)}]{\kn{ಪರಿಗಣಿಸು} (parigaṇisu) — to consider}

\label{lesson:KA-C381-pariganisu}



\begin{quote}
Before the new one: say the Kannada for an inducement, then the Kannada for to attract.
\end{quote}

\subsection*{You'll want to know: \kn{ಪರಿಗಣಿಸು}}
\textbf{\kn{ಪರಿಗಣಿಸು}} — \emph{parigaṇisu} — \textquotedblleft{}to consider\textquotedblright{}.

\begin{quote}
\kn{ನಾನು} \kn{ನಿನ್ನ} \kn{ಸಲಹೆಯನ್ನು} \kn{ಪರಿಗಣಿಸಿದೆ}.
\end{quote}

\begin{quote}
\emph{nānu ninna salaheyannu parigaṇiside.}
\end{quote}

\begin{quote}
I considered your suggestion.
\end{quote}

Past (I …): \kn{ಪರಿಗಣಿಸಿದೆ}, \emph{parigaṇiside}.

Also to count as: \kn{ಇದನ್ನು} \kn{ತಪ್ಪು} \kn{ಎಂದು} \kn{ಪರಿಗಣಿಸು}, to count this as a mistake.

\begin{cousinweb}
\textbf{\kn{ಪರಿಗಣಿಸು}} (\emph{parigaṇisu}) is built on Sanskrit \emph{pari-gaṇ}, \textquotedblleft{}to count over, to reckon up completely\textquotedblright{} (Monier-Williams), with the Kannada verb ending -\emph{isu}.
\end{cousinweb}

\begin{grammarlens}[title={What you've built}]
One more word for giving reasons and opinions.
\end{grammarlens}

\subsection*{Your turn}
\begin{itemize}
\raggedright
  \item \emph{Say it:} \emph{parigaṇisu}
  \item \emph{Say it:} \emph{parigaṇisu}, once more
  \item \emph{Say it:} \emph{nānu ninna salaheyannu parigaṇiside.} — I considered your suggestion.
  \item \emph{Recall:} say the Kannada for an inducement, then the Kannada for to attract, then say \emph{parigaṇisu} again
\end{itemize}

\begin{tcolorbox}[breakable,colback=teal!4,colframe=teal!35!black,title={Before you move on}]
What does \kn{ಪರಿಗಣಿಸು} mean? (\textquotedblleft{}To consider\textquotedblright{}.) Say it once more.
\end{tcolorbox}
\section[\texorpdfstring{\kn{ಭಾವಿಸು} (bhāvisu)}{ಭಾವಿಸು (bhāvisu)}]{\kn{ಭಾವಿಸು} (bhāvisu) — to suppose; to assume}

\label{lesson:KA-C381-bhavisu}



\begin{quote}
Before the new one: say the Kannada for to attract, then the Kannada for to consider.
\end{quote}

\subsection*{You'll want to know: \kn{ಭಾವಿಸು}}
\textbf{\kn{ಭಾವಿಸು}} — \emph{bhāvisu} — \textquotedblleft{}to suppose; to assume\textquotedblright{}.

\begin{quote}
\kn{ನೀನು} \kn{ಬರುವುದಿಲ್ಲ} \kn{ಅಂತ} \kn{ನಾನು} \kn{ಭಾವಿಸಿದೆ}.
\end{quote}

\begin{quote}
\emph{nīnu baruvudilla anta nānu bhāviside.}
\end{quote}

\begin{quote}
I assumed you would not come.
\end{quote}

Past (I …): \kn{ಭಾವಿಸಿದೆ}, \emph{bhāviside}.

What is supposed comes before \kn{ಅಂತ} or \kn{ಎಂದು}.

\begin{cousinweb}
\textbf{\kn{ಭಾವಿಸು}} (\emph{bhāvisu}) is built on Sanskrit \emph{bhāva}, \textquotedblleft{}being, state, disposition\textquotedblright{} (CDIAL 9475), from \emph{bhū}, \textquotedblleft{}to be\textquotedblright{}, with the Kannada verb ending -\emph{isu}.
\end{cousinweb}

\begin{grammarlens}[title={What you've built}]
One more word for giving reasons and opinions.
\end{grammarlens}

\subsection*{Your turn}
\begin{itemize}
\raggedright
  \item \emph{Say it:} \emph{bhāvisu}
  \item \emph{Say it:} \emph{bhāvisu}, once more
  \item \emph{Say it:} \emph{nīnu baruvudilla anta nānu bhāviside.} — I assumed you would not come.
  \item \emph{Recall:} say the Kannada for to attract, then the Kannada for to consider, then say \emph{bhāvisu} again
\end{itemize}

\begin{tcolorbox}[breakable,colback=teal!4,colframe=teal!35!black,title={Before you move on}]
What does \kn{ಭಾವಿಸು} mean? (\textquotedblleft{}To suppose; to assume\textquotedblright{}.) Say it once more.
\end{tcolorbox}
\section[\texorpdfstring{\kn{ವಿವೇಕ} (vivēka)}{ವಿವೇಕ (vivēka)}]{\kn{ವಿವೇಕ} (vivēka) — good sense; discernment}

\label{lesson:KA-C381-viveka}



\begin{quote}
Before the new one: say the Kannada for to consider, then the Kannada for to suppose.
\end{quote}

\subsection*{You'll want to know: \kn{ವಿವೇಕ}}
\textbf{\kn{ವಿವೇಕ}} — \emph{vivēka} — \textquotedblleft{}good sense; discernment\textquotedblright{}.

\begin{quote}
\kn{ಕಷ್ಟದ} \kn{ಸಮಯದಲ್ಲಿ} \kn{ಅವಳು} \kn{ವಿವೇಕದಿಂದ} \kn{ನಡೆದುಕೊಂಡಳು}.
\end{quote}

\begin{quote}
\emph{kaṣṭada samayadalli avaḷu vivēkadinda naḍedukoṇḍaḷu.}
\end{quote}

\begin{quote}
In hard times she acted with good sense.
\end{quote}

\kn{ವಿವೇಕಿ}, a sensible person.

\begin{cousinweb}
\textbf{\kn{ವಿವೇಕ}} (\emph{vivēka}) is borrowed from Sanskrit \emph{viveka}, \textquotedblleft{}discrimination, the power to tell things apart\textquotedblright{} (Monier-Williams).
\end{cousinweb}

\begin{grammarlens}[title={What you've built}]
One more word for giving reasons and opinions.
\end{grammarlens}

\subsection*{Your turn}
\begin{itemize}
\raggedright
  \item \emph{Say it:} \emph{vivēka}
  \item \emph{Say it:} \emph{vivēka}, once more
  \item \emph{Say it:} \emph{kaṣṭada samayadalli avaḷu vivēkadinda naḍedukoṇḍaḷu.} — In hard times she acted with good sense.
  \item \emph{Recall:} say the Kannada for to consider, then the Kannada for to suppose, then say \emph{vivēka} again
\end{itemize}

\begin{tcolorbox}[breakable,colback=teal!4,colframe=teal!35!black,title={Before you move on}]
What does \kn{ವಿವೇಕ} mean? (\textquotedblleft{}Good sense; discernment\textquotedblright{}.) Say it once more.
\end{tcolorbox}
\section[\texorpdfstring{\kn{ಮನನ} \kn{ಮಾಡು} (manana māḍu)}{ಮನನ ಮಾಡು (manana māḍu)}]{\kn{ಮನನ} \kn{ಮಾಡು} (manana māḍu) — to reflect on}

\label{lesson:KA-C381-manana-madu}



\begin{quote}
Before the new one: say the Kannada for to suppose, then the Kannada for good sense.
\end{quote}

\subsection*{You'll want to know: \kn{ಮನನ} \kn{ಮಾಡು}}
\textbf{\kn{ಮನನ} \kn{ಮಾಡು}} — \emph{manana māḍu} — \textquotedblleft{}to reflect on\textquotedblright{}.

\begin{quote}
\kn{ಗುರುಗಳ} \kn{ಮಾತನ್ನು} \kn{ನಾನು} \kn{ಮನೆಯಲ್ಲಿ} \kn{ಮನನ} \kn{ಮಾಡಿದೆ}.
\end{quote}

\begin{quote}
\emph{gurugaḷa mātannu nānu maneyalli manana māḍide.}
\end{quote}

\begin{quote}
At home I reflected on my teacher's words.
\end{quote}

Past (I …): \kn{ಮನನ} \kn{ಮಾಡಿದೆ}, \emph{manana māḍide}.

What you reflect on takes -\kn{ಅನ್ನು}.

\begin{cousinweb}
\textbf{\kn{ಮನನ} \kn{ಮಾಡು}} (\emph{manana māḍu}) is built from Sanskrit \emph{manana}, \textquotedblleft{}thinking, reflection\textquotedblright{} (Monier-Williams), from \emph{man}, \textquotedblleft{}to think\textquotedblright{}, with \emph{māḍu}, \textquotedblleft{}to do\textquotedblright{}.
\end{cousinweb}

\begin{grammarlens}[title={What you've built}]
One more word for giving reasons and opinions.
\end{grammarlens}

\subsection*{Your turn}
\begin{itemize}
\raggedright
  \item \emph{Say it:} \emph{manana māḍu}
  \item \emph{Say it:} \emph{manana māḍu}, once more
  \item \emph{Say it:} \emph{gurugaḷa mātannu nānu maneyalli manana māḍide.} — At home I reflected on my teacher's words.
  \item \emph{Recall:} say the Kannada for to suppose, then the Kannada for good sense, then say \emph{manana māḍu} again
\end{itemize}

\begin{tcolorbox}[breakable,colback=teal!4,colframe=teal!35!black,title={Before you move on}]
What does \kn{ಮನನ} \kn{ಮಾಡು} mean? (\textquotedblleft{}To reflect on\textquotedblright{}.) Say it once more.
\end{tcolorbox}
\section[\texorpdfstring{\kn{ಬುದ್ಧಿ} (buddhi)}{ಬುದ್ಧಿ (buddhi)}]{\kn{ಬುದ್ಧಿ} (buddhi) — intellect; wit}

\label{lesson:KA-C381-buddhi}



\begin{quote}
Before the new one: say the Kannada for good sense, then the Kannada for to reflect on.
\end{quote}

\subsection*{You'll want to know: \kn{ಬುದ್ಧಿ}}
\textbf{\kn{ಬುದ್ಧಿ}} — \emph{buddhi} — \textquotedblleft{}intellect; wit\textquotedblright{}.

\begin{quote}
\kn{ಆ} \kn{ಹುಡುಗನಿಗೆ} \kn{ಒಳ್ಳೆಯ} \kn{ಬುದ್ಧಿ} \kn{ಇದೆ}.
\end{quote}

\begin{quote}
\emph{ā huḍuganige oḷḷeya buddhi ide.}
\end{quote}

\begin{quote}
That boy has a good mind.
\end{quote}

Also advice: \kn{ಬುದ್ಧಿ} \kn{ಹೇಳು}, to give someone a piece of advice.

\begin{cousinweb}
\textbf{\kn{ಬುದ್ಧಿ}} (\emph{buddhi}) is borrowed from Sanskrit \emph{buddhi}, \textquotedblleft{}intelligence, understanding\textquotedblright{} (CDIAL 9277), from \emph{budh}, \textquotedblleft{}to wake, to know\textquotedblright{}.
\end{cousinweb}

\begin{grammarlens}[title={What you've built}]
One more word for giving reasons and opinions.
\end{grammarlens}

\subsection*{Your turn}
\begin{itemize}
\raggedright
  \item \emph{Say it:} \emph{buddhi}
  \item \emph{Say it:} \emph{buddhi}, once more
  \item \emph{Say it:} \emph{ā huḍuganige oḷḷeya buddhi ide.} — That boy has a good mind.
  \item \emph{Recall:} say the Kannada for good sense, then the Kannada for to reflect on, then say \emph{buddhi} again
\end{itemize}

\begin{tcolorbox}[breakable,colback=teal!4,colframe=teal!35!black,title={Before you move on}]
What does \kn{ಬುದ್ಧಿ} mean? (\textquotedblleft{}Intellect; wit\textquotedblright{}.) Say it once more.
\end{tcolorbox}
